\documentclass[aspectratio=169]{beamer}

% ==== Theme & basic look ====
\usetheme{Madrid}
\usecolortheme{seahorse}
\setbeamertemplate{navigation symbols}{}
\setbeamertemplate{footline}[frame number]


% Packages for mathematics  
\usepackage{amsmath}  
\usepackage{amsfonts}  
\usepackage{amssymb}  
\usepackage{CJKutf8}
% Support for \Tilde macro used throughout
\providecommand{\Tilde}[1]{\tilde{#1}}
\providecommand{\proofname}{Proof}
\newenvironment<>{proofs}[1][\proofname]{%
    \par
    \def\insertproofname{#1.}%
    \usebeamertemplate{proof begin}#2}
  {\usebeamertemplate{proof end}}
\makeatother

% Math shorthand and vector symbols
\newcommand{\grad}{\nabla}
\newcommand{\EE}{\mathbb{E}}
\newcommand{\dd}{\mathrm{d}}
\newcommand{\kbT}{k_B T}
\newcommand{\vx}{\mathbf{x}}

% Bold vector/time-indexed variables
\newcommand{\vxz}{\mathbf{x}_0}
\newcommand{\vxt}{\mathbf{x}_t}
\newcommand{\alpt}{\alpha_t}
\newcommand{\sigt}{\sigma_t}
\newcommand{\veps}{\boldsymbol{\epsilon}}
\newcommand{\vF}{\mathbf{F}}

% Discrete diffusion specific notation
\newcommand{\vq}{\mathbf{q}}
\newcommand{\vQ}{\mathbf{Q}}
\newcommand{\vK}{\mathbf{K}}
\newcommand{\vV}{\mathbf{V}}
\newcommand{\vP}{\mathbf{P}}

% Title page details:   
\title[Discrete diffusion models]{Discrete diffusion models: From continuous to categorical}  
\author{Pipi Hu}  
\institute[BIMSA]{Beijing Institute of Mathematical Sciences and Applications}  
\date{\today}  

\pgfdeclareimage[width=2cm]{logo}{../figures/logo.png}
\logo{\pgfuseimage{logo}}

% Outline slide (moved to later; avoid frame before \begin{document})

\setbeamertemplate{navigation symbols}{}
\AtBeginSection[]  
{  
    \begin{frame}<beamer>{Outline}         
        \tableofcontents[currentsection]  
    \end{frame}  
}  
  
\AtBeginSubsection[]  
{  
    \begin{frame}<beamer>{Outline}         
        \tableofcontents[currentsection,currentsubsection]  
    \end{frame}  
}  

\begin{document}  
  
% Title slide  
\begin{frame}  
  \titlepage  
\end{frame}  

\begin{frame}{Feynman's words}
\centering
    What I cannot create, I do not understand.
\end{frame}
  
% Presentation slides  

\section{From Continuous to Discrete: The Challenge}

\begin{frame}{The discrete data challenge}
Traditional diffusion models work well for continuous data, but many real-world applications involve discrete data:
\begin{itemize}
    \item \textcolor{red}{Text generation}: Words, tokens, characters
    \item \textcolor{red}{Image generation}: Pixel values, categorical features
    \item \textcolor{red}{Molecular generation}: Atom types, bond types
    \item \textcolor{red}{Music generation}: Notes, chords, instruments
\end{itemize}

\textbf{Key question}: How to apply diffusion principles to discrete/categorical data?
\end{frame}

\begin{frame}{Why not just discretize continuous diffusion?}
Direct discretization of continuous diffusion faces several challenges:
\begin{enumerate}
    \item \textcolor{red}{Score function undefined}: $\nabla \log p(x)$ doesn't exist for discrete $x$
    \item \textcolor{red}{Noise structure}: Gaussian noise doesn't preserve categorical structure
    \item \textcolor{red}{Sampling complexity}: Discrete spaces have different geometry
    \item \textcolor{red}{Training instability}: Continuous losses don't translate directly
\end{enumerate}

\textbf{Need}: New mathematical framework for discrete diffusion
\end{frame}

\section{Discrete Diffusion: Mathematical Foundation}

\begin{frame}{Discrete state spaces and transition matrices}
For discrete data $x \in \{1, 2, \ldots, K\}$, we need discrete transition probabilities.

\textbf{Forward process}: Define transition matrix $\vQ_t$ such that
\begin{equation}
    q(x_t | x_{t-1}) = \vQ_t[x_{t-1}, x_t]
\end{equation}

\textbf{Key properties}:
\begin{itemize}
    \item $\vQ_t$ is a stochastic matrix: $\sum_{x_t} \vQ_t[x_{t-1}, x_t] = 1$
    \item $\vQ_t[x_{t-1}, x_t] \geq 0$ for all $x_{t-1}, x_t$
    \item Can be time-dependent: $\vQ_t$ changes with $t$
\end{itemize}
\end{frame}

\begin{frame}{Discrete forward process}
The discrete forward process is defined as:
\begin{equation}
    q(x_t | x_0) = \sum_{x_1, \ldots, x_{t-1}} \prod_{s=1}^t \vQ_s[x_{s-1}, x_s]
\end{equation}

For \textcolor{red}{uniform transitions} (common choice):
\begin{equation}
    \vQ_t = (1-\beta_t)\mathbf{I} + \beta_t \mathbf{1}\mathbf{1}^T/K
\end{equation}
where $\mathbf{I}$ is identity matrix and $\mathbf{1}\mathbf{1}^T$ is uniform transition matrix.

\textbf{Interpretation}: With probability $(1-\beta_t)$, stay the same; with probability $\beta_t$, jump to any state uniformly.
\end{frame}

\begin{frame}{Discrete reverse process}
The reverse process requires learning transition probabilities:
\begin{equation}
    p_\theta(x_{t-1} | x_t) = \text{Cat}(x_{t-1}; \boldsymbol{\pi}_\theta(x_t, t))
\end{equation}

where $\boldsymbol{\pi}_\theta(x_t, t) \in \mathbb{R}^K$ is a probability vector.

\textbf{Training objective}: Learn to reverse the forward process
\begin{equation}
    \mathcal{L} = \mathbb{E}_{t, x_0, x_t} \left[ -\log p_\theta(x_{t-1} | x_t) \right]
\end{equation}

\textbf{Challenge}: How to parameterize $\boldsymbol{\pi}_\theta(x_t, t)$ effectively?
\end{frame}

\section{Score-based Discrete Diffusion}

\begin{frame}{Discrete score function}
For discrete data, we define the \textcolor{red}{discrete score function}:
\begin{equation}
    s(x, t) = \nabla \log p(x, t) = \frac{\partial \log p(x, t)}{\partial x}
\end{equation}

\textbf{Problem}: This is not well-defined for categorical $x$.

\textbf{Solution}: Use \textcolor{red}{logits} instead of direct gradients:
\begin{equation}
    s(x, t) = \log p(x, t) - \log \sum_{x'} p(x', t)
\end{equation}

Or equivalently:
\begin{equation}
    s(x, t) = \log p(x, t) + C
\end{equation}
where $C$ is a constant that doesn't affect the relative probabilities.
\end{frame}

\begin{frame}{Discrete denoising score matching}
For discrete data, the DSM loss becomes:
\begin{equation}
    \mathcal{L}_{DSM} = \mathbb{E}_{t, x_0, x_t} \left[ \|s_\theta(x_t, t) - \nabla_{x_t} \log q(x_t | x_0)\|^2 \right]
\end{equation}

\textbf{Key insight}: The score $\nabla_{x_t} \log q(x_t | x_0)$ can be computed analytically for discrete transitions.

For uniform transitions:
\begin{equation}
    \nabla_{x_t} \log q(x_t | x_0) = \frac{\vQ_t[x_0, x_t]}{\sum_{x'} \vQ_t[x_0, x']} - \frac{1}{K}
\end{equation}
\end{frame}

\section{Autoregressive vs Diffusion for Discrete Data}

\begin{frame}{Autoregressive models for discrete data}
Traditional approach for discrete data:
\begin{equation}
    p(x) = \prod_{i=1}^n p(x_i | x_{<i})
\end{equation}

\textbf{Advantages}:
\begin{itemize}
    \item Natural for discrete data
    \item Easy to train with cross-entropy loss
    \item Can model complex dependencies
\end{itemize}

\textbf{Disadvantages}:
\begin{itemize}
    \item Sequential generation (slow)
    \item Order dependence
    \item Hard to parallelize
\end{itemize}
\end{frame}

\begin{frame}{Discrete diffusion advantages}
Discrete diffusion offers several advantages over autoregressive models:
\begin{enumerate}
    \item \textcolor{red}{Parallel generation}: All positions can be generated simultaneously
    \item \textcolor{red}{Order independence}: No fixed generation order
    \item \textcolor{red}{Flexible conditioning}: Easy to condition on partial information
    \item \textcolor{red}{Unified framework}: Same principles as continuous diffusion
\end{enumerate}

\textbf{Trade-off}: More complex training but more flexible generation
\end{frame}

\section{Implementation Strategies}

\begin{frame}{Discrete diffusion implementation}
\textbf{Strategy 1: Direct parameterization}
\begin{equation}
    p_\theta(x_{t-1} | x_t) = \text{softmax}(f_\theta(x_t, t))
\end{equation}
where $f_\theta: \mathbb{R}^K \times \mathbb{R} \to \mathbb{R}^K$

\textbf{Strategy 2: Score-based approach}
\begin{equation}
    s_\theta(x_t, t) = f_\theta(x_t, t) - \log \sum_{x'} \exp(f_\theta(x', t))
\end{equation}

\textbf{Strategy 3: Embedding approach}
\begin{itemize}
    \item Embed discrete states in continuous space
    \item Apply continuous diffusion in embedding space
    \item Project back to discrete space
\end{itemize}
\end{frame}

\begin{frame}{Training discrete diffusion models}
\textbf{Training procedure}:
\begin{enumerate}
    \item Sample $x_0 \sim p_{data}$
    \item Sample $t \sim \text{Uniform}(0, T)$
    \item Sample $x_t \sim q(x_t | x_0)$ using forward process
    \item Train model to predict $x_0$ or $x_{t-1}$ from $x_t$
\end{enumerate}

\textbf{Loss functions}:
\begin{itemize}
    \item Cross-entropy: $\mathcal{L} = -\log p_\theta(x_0 | x_t)$
    \item Score matching: $\mathcal{L} = \|s_\theta(x_t, t) - s_{true}(x_t, t)\|^2$
    \item Variational bound: $\mathcal{L} = \text{ELBO}(x_0, x_t)$
\end{itemize}
\end{frame}

\section{Advanced Discrete Diffusion}

\begin{frame}{Categorical diffusion}
For categorical data with $K$ classes, use categorical distribution:
\begin{equation}
    q(x_t | x_{t-1}) = \text{Cat}(x_t; \boldsymbol{\alpha}_t)
\end{equation}

where $\boldsymbol{\alpha}_t$ is the concentration parameter.

\textbf{Forward process}:
\begin{equation}
    \boldsymbol{\alpha}_t = (1-\beta_t) \mathbf{e}_{x_{t-1}} + \beta_t \mathbf{1}/K
\end{equation}

where $\mathbf{e}_{x_{t-1}}$ is the one-hot vector for $x_{t-1}$.
\end{frame}

\begin{frame}{Discrete diffusion with structure}
For structured discrete data (e.g., molecules, graphs):
\begin{itemize}
    \item \textcolor{red}{Graph diffusion}: Diffuse over graph structure
    \item \textcolor{red}{Hierarchical diffusion}: Multi-level discrete states
    \item \textcolor{red}{Conditional diffusion}: Condition on partial structure
\end{itemize}

\textbf{Example}: Molecular generation
\begin{equation}
    p(\text{molecule}) = \prod_{t=1}^T p(\text{atom}_t | \text{atoms}_{<t}, \text{bonds}_{<t})
\end{equation}

\textbf{Challenge}: Maintaining chemical validity during diffusion
\end{frame}

\section{Applications and Examples}

\begin{frame}{Text generation with discrete diffusion}
\textbf{Token-level diffusion}:
\begin{itemize}
    \item Forward: Gradually replace tokens with noise
    \item Reverse: Learn to denoise and recover original text
    \item Training: Predict original tokens from noisy versions
\end{itemize}

\textbf{Advantages over autoregressive}:
\begin{itemize}
    \item Parallel generation
    \item Better long-range dependencies
    \item More flexible conditioning
\end{itemize}
\end{frame}

\begin{frame}{Image generation with discrete diffusion}
\textbf{Pixel-level discrete diffusion}:
\begin{itemize}
    \item Treat pixel values as discrete categories
    \item Forward: Gradually randomize pixel values
    \item Reverse: Learn to recover original pixels
\end{itemize}

\textbf{Hybrid approaches}:
\begin{itemize}
    \item Continuous diffusion for low-level features
    \item Discrete diffusion for high-level structure
    \item Multi-scale discrete diffusion
\end{itemize}
\end{frame}

\section{Theoretical Analysis}

\begin{frame}{Discrete diffusion convergence}
\textbf{Theorem}: Under certain conditions, discrete diffusion converges to the data distribution.

\textbf{Key assumptions}:
\begin{enumerate}
    \item Forward process is ergodic
    \item Reverse process is properly parameterized
    \item Training loss is minimized
\end{enumerate}

\textbf{Proof sketch}:
\begin{equation}
    p_\theta(x_0) = \int p_\theta(x_0 | x_T) p(x_T) dx_T
\end{equation}

As $T \to \infty$, $p(x_T)$ approaches uniform distribution, and $p_\theta(x_0 | x_T)$ learns to recover the data distribution.
\end{frame}

\begin{frame}{Discrete vs continuous diffusion}
\textbf{Comparison}:
\begin{itemize}
    \item \textcolor{red}{Continuous}: Smooth interpolation, gradient-based
    \item \textcolor{red}{Discrete}: Categorical jumps, probability-based
\end{itemize}

\textbf{Mathematical differences}:
\begin{itemize}
    \item Score function: $\nabla \log p(x)$ vs $\log p(x)$
    \item Noise: Gaussian vs categorical
    \item Sampling: Langevin dynamics vs discrete transitions
\end{itemize}

\textbf{Unified view}: Both can be seen as special cases of a general diffusion framework
\end{frame}

\section{Challenges and Open Problems}

\begin{frame}{Current challenges in discrete diffusion}
\begin{enumerate}
    \item \textcolor{red}{Training stability}: Discrete losses can be unstable
    \item \textcolor{red}{Sampling efficiency}: Discrete sampling can be slow
    \item \textcolor{red}{Model capacity}: Need large models for complex discrete data
    \item \textcolor{red}{Evaluation metrics}: How to measure quality of discrete generation?
\end{enumerate}

\textbf{Open problems}:
\begin{itemize}
    \item Optimal transition matrices for different data types
    \item Theoretical guarantees for discrete diffusion
    \item Efficient sampling algorithms
    \item Better loss functions for discrete data
\end{itemize}
\end{frame}

\begin{frame}{Future directions}
\textbf{Research directions}:
\begin{enumerate}
    \item \textcolor{red}{Hybrid models}: Combine continuous and discrete diffusion
    \item \textcolor{red}{Structured diffusion}: Incorporate domain knowledge
    \item \textcolor{red}{Efficient algorithms}: Faster training and sampling
    \item \textcolor{red}{Theoretical understanding}: Better mathematical foundations
\end{enumerate}

\textbf{Applications}:
\begin{itemize}
    \item Scientific discovery (molecules, materials)
    \item Creative applications (music, art)
    \item Natural language processing
    \item Code generation
\end{itemize}
\end{frame}

\section{Conclusion}

\begin{frame}{Summary}
Discrete diffusion models provide a powerful framework for generating discrete data:
\begin{itemize}
    \item \textcolor{red}{Mathematical foundation}: Discrete transition matrices and score functions
    \item \textcolor{red}{Implementation strategies}: Direct parameterization, score-based, embedding-based
    \item \textcolor{red}{Applications}: Text, images, molecules, and more
    \item \textcolor{red}{Advantages}: Parallel generation, flexible conditioning, unified framework
\end{itemize}

\textbf{Key insight}: Discrete diffusion extends the power of continuous diffusion to categorical data while maintaining the core principles of iterative denoising.
\end{frame}

\begin{frame}{Questions}
If you are interested in the questions, please feel free to write an email to me {\color{red}(\href{mailto:hpp@bimsa.cn}{hpp@bimsa.cn})} with your comments. 
\begin{enumerate}[(1.)]
    \item How can we design optimal transition matrices for specific types of discrete data (e.g., text, molecules, graphs)?
    \item What are the theoretical guarantees for convergence of discrete diffusion models?
    \item How can we efficiently sample from discrete diffusion models while maintaining quality?
    \item What are the best practices for training discrete diffusion models on large-scale datasets?
\end{enumerate}
\end{frame}

% Thank you slide  
\begin{frame}{Welcome inputs}  
  \centering  
  Any comments?
  \\[2em]
  Welcome your inputs!  
\end{frame}  
  
\end{document}  
